% !TEX program = xelatex
% Paper-talk Beamer skeleton — 16:9, back-row-readable defaults.
% Works for English decks out of the box; for Chinese decks uncomment the
% ctex block. Compile with XeLaTeX.
\documentclass[aspectratio=169,12pt]{beamer}

% ---- theme: pick one sober theme; replace with your institutional one ----
\usetheme{Madrid}
\usecolortheme{seahorse}
\setbeamertemplate{navigation symbols}{}   % navigation icons are noise
\setbeamertemplate{footline}[frame number]

% ---- fonts: Latin/numbers in Times; CJK (if any) in a hei-style face ----
\usepackage{newtxtext}
% \usepackage{ctex}
% \setCJKmainfont[BoldFont={Microsoft YaHei Bold}]{Microsoft YaHei}
% \setCJKsansfont[BoldFont={Microsoft YaHei Bold}]{Microsoft YaHei}

\usepackage{graphicx,booktabs}

% ---- two small helpers every deck ends up wanting ----
% \hl      highlight the one number that matters, in the theme color
% \srcnote tiny gray source line under figures/tables
\newcommand{\hl}[1]{{\color{structure.fg}\bfseries #1}}
\newcommand{\srcnote}[1]{{\tiny\color{gray} #1}}

\title[Short Title]{Full Title of the Talk}
\subtitle{A Paper Presentation}
\author[Author]{Author Name}
\institute[Institute Short]{Institute Full Name}
\date{\today}

\begin{document}

\begin{frame}
  \titlepage
\end{frame}

\begin{frame}{Outline}
  \tableofcontents
\end{frame}

\section{Background}

% ---- pattern 1: content page — the load-bearing sentences, full width ----
\begin{frame}{One Idea Per Page}
  \begin{itemize}
    \setlength{\itemsep}{6pt}
    \item The single most important sentence from this stretch of the script,
          with its key number \hl{highlighted}.
    \item At most five or six short items. If you need more, this is two pages.
    \item Enough substance that the speaker can glance at the screen and
          remember what comes next — the deck is the speaker's prompter too.
  \end{itemize}
  \vfill
  \srcnote{Source: where the numbers on this page come from.}
\end{frame}

% ---- pattern 2: figure page — big figure + secondary script text beside ----
% The figure is what the audience studies; the side text is what the speaker
% reads aloud meanwhile (mark it :::verbatim in the script).
\begin{frame}{Figure Page}
  \begin{columns}[T]
    \column{0.42\textwidth}
      {\small Secondary material: context and caveats that didn't make the
       key-points cut. Smaller font is fine here — it is the speaker's text,
       shown on screen while the audience looks at the figure.\par}
    \column{0.55\textwidth}
      \centering
      % \includegraphics[width=\linewidth]{figure}
      \fbox{\parbox[c][0.55\textheight][c]{0.9\linewidth}{\centering figure}}
      \\[2pt]
      \srcnote{What the figure shows, and where it is from.}
  \end{columns}
\end{frame}

\section{Results}
% ...

\end{document}
